\section{Coherent and preparation directions at a singular boundary}\label{sec:coherent70}
The preparation theorem has a square-root horizon scale. Reversible
quantum information can instead accumulate a coherent discrepancy at a
linear scale. We now exhibit both scales in a single non-input-erasing
family, and determine their combined description cost.

Fix $d\ge2$, $n,m\ge1$ and public rank bounds $0\le r_y\le n$, not all
zero. For $U\in U(d)$ and $\sigma\in\mathcal S_{\boldsymbol r}$ define
an instrument from $M_d$ to $M_{dn}$ by
\begin{equation}\label{eq:coherentfamily70}
 \mathcal F_{U,\sigma,y}(X)=UXU^*\otimes\sigma_y.
\end{equation}
Its outcome register is retained. Global phase is immaterial, so the
unitary parameter lies in $PU(d)$. Put
\[
 u=d^2-1,\qquad v=\sum_y r_y(2n-r_y)-1.
\]
The Choi matrix of each outcome is, up to a fixed tensor permutation,
$|U\rangle\!\rangle\langle\!\langle U|\otimes\sigma_y$. Its rank is
$\operatorname{rank}\sigma_y$, not $d\operatorname{rank}\sigma_y$ as in
the input-erasing family. These instruments transmit the data system
unitarily, including its entanglement with a reference. Their preparation
supports may change and may lie on any of the permitted singular strata.
Nevertheless, the outcome probabilities do not depend on the input. This
is a specified boundary family, not the whole body of disturbing
instruments.

Let $\mathcal C_N(\delta;\boldsymbol r)$ denote its covering number in
$d_N$, with centres allowed to be any legal memoryless instrument on the
same input, output and outcome spaces. Let $\mathcal C_N^{\Q}$ impose the
additional condition that the centres have Gaussian-rational Choi entries.
A target is supplied as $U$ and its preparation blocks; no unknown-device
learning interface is used.

\begin{theorem}[Joint coding at a coherent--preparation boundary]\label{thm:coherententropy70}
There are $0<c\le C<\infty$, depending only on the fixed dimensions and
rank bounds, such that, for every $N\ge1$ and $0<\delta\le1/32$,
\begin{equation}\label{eq:coherententropy70}
 c\left(\frac N\delta\right)^u
  \left(\frac{\sqrt N}\delta\right)^v
 \le\mathcal C_N(\delta;\boldsymbol r)
 \le\mathcal C_N^{\Q}(\delta;\boldsymbol r)
 \le C\left(\frac N\delta\right)^u
  \left(\frac{\sqrt N}\delta\right)^v.
\end{equation}
The metric is the unhalved final-state trace norm, maximized over common
adaptive quantum testers using the same memoryless instrument at most
$N$ times, with arbitrary finite references, quantum memory, feedback and
bounded public stopping. In particular, the optimal fixed-length reusable
description has length
\begin{equation}\label{eq:coherentbits70}
 (u+v/2)\log_2N+(u+v)\log_2(1/\delta)+O_{d,n,m,\boldsymbol r}(1).
\end{equation}
The upper centres can be chosen within the family
\eqref{eq:coherentfamily70}, with rational unitary matrices and rational
preparation blocks. They do not increase a target block rank and preserve
its zero outcomes. A finite exact encoder exists for supplied rational
unitaries and rational preparation blocks. Description length is not its
workspace or running time, and decoding describes a quantum instrument,
not a classical physical processor for unknown entangled inputs.
\end{theorem}

\subsection{The two metric scales}
For $U,V\in U(d)$ write
\begin{equation}\label{eq:chordunitary70}
 \chi(U,V)=\left(1-\frac{|\tr(U^*V)|^2}{d^2}\right)^{1/2},\qquad
 \Delta_N(\sigma,\tau)=
       \|\boldsymbol\sigma^{\otimes N}-\boldsymbol\tau^{\otimes N}\|_1.
\end{equation}
The first expression is the chordal distance between normalized Choi rays;
it is a metric after quotienting global phase.

\begin{lemma}[Anisotropic adaptive comparison]\label{lem:coherentmetric70}
For the instruments \eqref{eq:coherentfamily70},
\begin{align}
 \max\{\min\{N\chi(U,V)/\sqrt2,1\},\Delta_N(\sigma,\tau)\}
  &\le d_N(\mathcal F_{U,\sigma},\mathcal F_{V,\tau}),\label{eq:coherentlower70}\\
 d_N(\mathcal F_{U,\sigma},\mathcal F_{V,\tau})
  &\le\min\{2,2dN\chi(U,V)+\Delta_N(\sigma,\tau)\}.
                                                        \label{eq:coherentupper70}
\end{align}
All lower testers are permitted to depend on the pair being compared.
\end{lemma}
\begin{proof}
Discarding the data-system outputs gives the product-state lower bound.
For the unitary bound, discard the preparation systems and outcome labels,
and apply the same known operation $U^*$ to the data output after each
call. The two resulting data operations are $I$ and $W=U^*V$.
If $e^{i\phi_j}$ are the eigenvalues of $W$, then
\[
 \chi(U,V)^2=\frac2{d^2}\sum_{i,j}
                  \sin^2((\phi_i-\phi_j)/2).
\]
There is an eigenvalue pair whose principal half-angle
$\theta\in[0,\pi/2]$ satisfies $\sin\theta\ge\chi(U,V)/\sqrt2$.
Prepare the equal superposition of its orthogonal eigenvectors. After
$k$ calls its two possible pure states have trace distance
$2|\sin(k\theta)|$. If $N\theta\le\pi/2$, take $k=N$ and use
$\sin x\ge2x/\pi$ on $[0,\pi/2]$. The result is at least
$N\sin\theta$. If $N\theta>\pi/2$, take
$k=\lceil\pi/(4\theta)\rceil\le N$. Then
$\pi/4\le k\theta\le3\pi/4$, giving distance at least $\sqrt2$.
The case $\theta=0$ is immediate. This proves the first minimum in
\eqref{eq:coherentlower70}; the two lower witnesses need not coincide.

For the upper bound first change $U$ to $V$, keeping $\sigma$ fixed.
The unnormalized Choi difference of the unitary channels has trace norm
$2d\chi(U,V)$ and bounds their diamond distance. Appending a fixed
preparation state does not increase that norm. A slot-by-slot hybrid and
trace contraction for the common tester therefore cost at most
$2dN\chi(U,V)$. Next hold $V$ fixed and change $\sigma$ to $\tau$.
All calls to $V$ and all tester operations form one common quantum
processor consuming $N$ preparation states; unused copies are discarded
at a public stop. This costs at most $\Delta_N(\sigma,\tau)$, not the
sum of $N$ one-copy trace distances. Triangle inequality gives
\eqref{eq:coherentupper70}.
\end{proof}

Finite-use amplification of unitary differences is classical; see
Ac\'in \cite{Acin70} and Duan--Feng--Ying \cite{DFY70}. The elementary
estimate above is included to specify its uniform constant, its
pair-dependent controls and its compatibility with bounded stopping.
Neither the existence of unitary discrimination nor the real dimension
$d^2-1$ is claimed as a new ingredient.

\subsection{A finite rational atlas for unitary channels}
For $\ell\ge1$, a signed stereographic chart on $S^\ell$ has the form
\begin{equation}\label{eq:sphereatlas70}
 S_{j,\xi}(t)_j=\xi\frac{1-\|t\|^2}{1+\|t\|^2},\qquad
 S_{j,\xi}(t)_i=\xi\frac{2t_i}{1+\|t\|^2}\quad(i\ne j),
 \qquad \xi\in\{-1,1\}.
\end{equation}
For a target $q\in S^\ell$, choose $j$ of largest absolute coordinate,
$\xi=\operatorname{sgn}(q_j)$ and
$t_i=\xi q_i/(1+|q_j|)$. Then $|t_i|\le1$ and
$S_{j,\xi}(t)=q$. Direct expansion gives
\begin{equation}\label{eq:stereodistance70}
 \|S_{j,\xi}(t)-S_{j,\xi}(t')\|
 =\frac{2\|t-t'\|}{\sqrt{(1+\|t\|^2)(1+\|t'\|^2)}}.
\end{equation}
Rounding each coordinate towards zero at denominator $B$ thus gives a
rational unit vector at distance at most $2\sqrt\ell/B$. Every decoded
chart word is on the sphere, regardless of whether it is the canonical
encoding of a target.

\begin{lemma}[Rational unitary code]\label{lem:unitaryatlas70}
Put $L=d(d-1)/2$ and $K_d=3L+2(d-1)$. There is an explicitly enumerable
set $\mathcal U_B\subset U(d,\Q(i))$ satisfying
\begin{equation}\label{eq:unitarycode70}
 |\mathcal U_B|\le6^L4^{d-1}(2B+1)^{d^2-1},\qquad
 \sup_U\min_{V\in\mathcal U_B}\chi(U,V)\le K_d/B.
\end{equation}
For rational $U$, membership in any prescribed rational squared-error
ball $\chi(U,V)^2\le b$ is decidable by rational arithmetic alone.
\end{lemma}
\begin{proof}
The usual complex Givens elimination, in a fixed order of index pairs,
expresses every unitary as a product of $L$ embedded matrices
\[
 \begin{pmatrix}c&-\overline s\\s&c\end{pmatrix},
       \qquad c\in\R,\quad c^2+|s|^2=1,
\]
and a diagonal unitary. For completeness, a column pair $(a,b)$ is sent
to one nonzero coordinate using $c=|a|/(|a|^2+|b|^2)^{1/2}$ and
$s=cb/a$ when $a\ne0$; the cases $a=0$ or $(a,b)=0$ are obtained by
$c=0$ or by the identity. Applying the adjoint rotation eliminates the
second coordinate. Successive eliminations give an upper triangular
unitary, necessarily diagonal. Removing the phase of its first diagonal
entry leaves $d-1$ free phases. All factors are taken in the corresponding
fixed reversed order when reconstructing $U$.

Each Givens matrix has an $S^2$ parameter $(c,\Re s,\Im s)$; its
operator-norm change equals the Euclidean change of that vector. Each
remaining phase has an $S^1$ parameter. Apply
\eqref{eq:sphereatlas70} at grid $B$ to these $L+d-1$ factors.
There are $6$ chart headers for each $S^2$ and $4$ for each $S^1$,
with $2L+d-1=d^2-1$ grid digits altogether. Telescoping a product of
unitaries gives, up to a global phase, operator error at most
\[
 (2\sqrt2 L+2(d-1))/B\le K_d/B.
\]
For any fixed phase $z$ of modulus one,
$\chi(U,V)^2\le d^{-1}\|U-zV\|_F^2\le\|U-zV\|_{\rm op}^2$.
This proves \eqref{eq:unitarycode70}. The formula
\eqref{eq:chordunitary70} has rational square for rational $U,V$, proving
the last assertion. Enumeration of these finite words followed by the
squared-distance test is an exact encoder whenever the radius is at least
$K_d/B$. This does not assert a fast nearest-centre algorithm.
\end{proof}

\subsection{The joint covering theorem}
\begin{proof}[Proof of Theorem~\ref{thm:coherententropy70}]
For the upper bound use
\[
 B_U=\left\lceil4dK_dN/\delta\right\rceil,
 \qquad B_P=\max\{1,\lceil8\sqrt{Nv}/\delta\rceil\}.
\]
Lemma~\ref{lem:unitaryatlas70} makes the first term in
\eqref{eq:coherentupper70} at most $\delta/2$.
The factor code in Theorem~\ref{thm:factorcode68} makes the second term
at most $\delta/2$. Its finite pivot headers and anchor code give at most
$C_{n,m,\boldsymbol r}(2B_P+1)^v$ legal images. If $v=0$ the preparation
family is a singleton and no such digits are needed. Multiplication of
the two cardinalities proves the upper bound and all legality, rank and
zero-outcome assertions. For rational data, search the finite unitary
list using the rational test
$\chi(U,V)^2\le\delta^2/(16d^2N^2)$; the preceding estimate guarantees
a result. The preparation encoder is already exact. Decoded Choi entries
are rational products of the two rational outputs. The finite search and
the expanded matrices are not charged to the reusable payload.

For the converse, a neighborhood of the identity in $PU(d)$ contains a
quantitatively bi-Lipschitz $u$-dimensional Euclidean patch for the metric
$\chi$. Here is a direct verification. Use traceless Hermitian coordinates
$H$ near zero and the ray of $e^{iH}$. The differential of its normalized
Choi projection at zero vanishes only when $H$ is scalar, hence only when
$H=0$. On a sufficiently small convex coordinate ball its differential
is uniformly closer than half its least singular value to that at zero.
Integration along a line segment gives both Lipschitz inequalities.
Since the projection Frobenius distance is $\sqrt2\chi$, the asserted
metric comparison follows. Volumetric packing therefore gives at least
$c_d(N/\delta)^u$ unitaries separated by more than $4\delta/N$, with
a smaller fixed constant covering the coarse-scale cases as well.

The maximal-rank preparation patch used in
Theorem~\ref{thm:rankentropy68} has real dimension $v$ and is
bi-Lipschitz in tuple Frobenius norm. Pack it at separation greater than
$32\delta/\sqrt N$, obtaining at least
$c_{n,m,\boldsymbol r}(\sqrt N/\delta)^v$ points. The local binary
bound in Corollary~\ref{cor:localchoi68}, applied to the preparation
state (input dimension one), separates their $N$-copy states by more
than $2\delta$. Indeed the bound at that scale is
$(3/16)\min\{16\delta,1\}$, which is $3\delta$ before saturation
and at least $3/16>2\delta$ afterwards. For $v=0$ use its single point.

Take the Cartesian product of these packings. Two elements with different
unitaries are separated by \eqref{eq:coherentlower70}, since
$\min\{2\sqrt2\delta,1\}>2\delta$. Two elements with the same unitary
and different preparation blocks are separated by the other lower witness.
Thus a ball of radius $\delta$ about any legal centre can contain at most
one product point. This proves the lower bound even for centres outside
the family. Taking base-two logarithms and rounding a fixed-length word
length upward changes the answer by at most a fixed constant, proving
\eqref{eq:coherentbits70}.
\end{proof}

The coefficient $u+v/2$ reflects two different reuse scales, rather than
half of the total parameter dimension $u+v$. No spectral-gap assumption
or positive eigenvalue margin is used. The constants are nevertheless
fixed-dimensional. Nor is this a classification of every tangent direction
of an arbitrary instrument: coupling an input-dependent measurement to
its disturbance is not part of \eqref{eq:coherentfamily70}.

\subsection{An explicit qubit codec}
For $q=(q_0,q_1,q_2,q_3)\in S^3$ put
\begin{equation}\label{eq:quaternion70}
 U(q)=\begin{pmatrix}q_0+iq_3&q_2+iq_1\\-q_2+iq_1&q_0-iq_3\end{pmatrix}.
\end{equation}
The points $q$ and $-q$ describe the same channel. Choose a largest
absolute coordinate, flip that common sign to make it positive, and use
\eqref{eq:sphereatlas70} with $\ell=3$ and $\xi=1$. Its header has
four possibilities, and its body has exactly three signed digits.

\begin{corollary}[Exact rational coherent--preparation description]\label{cor:qubitcodec70}
For a supplied rational unit quaternion and rational preparation blocks,
the preceding three-digit chart, at $B_U=\lceil16N/\delta\rceil$, combined
with the preparation codec at tolerance $\delta/2$, gives a legal rational
instrument within $d_N$ error $\delta$. Its fixed-length description has
\[
 (3+v/2)\log_2N+(3+v)\log_2(1/\delta)+O_{n,m,\boldsymbol r}(1)
\]
bits, with no increase of a block rank and exact preservation of zero
outcomes. It uses only integer and rational arithmetic, including the
integer square roots in the preparation code. Its output is a mathematical
instrument description, not physical classical simulation.
\end{corollary}
\begin{proof}
For quaternion matrices,
$\|U(q)-U(q')\|_{\rm op}=\|q-q'\|$. The unitary diamond error is
at most twice this quantity. Thus the chart's coherent contribution
is at most $4\sqrt3 N/B_U<\delta/2$. Its squared upper certificate is
$48N^2/B_U^2$. The preparation contribution is at most $\delta/2$.
Triangle inequality proves the claim; one may certify its square by
$2(e_U^2+e_P^2)$, not by $e_U^2+e_P^2$. The chart returns an exactly
unit quaternion, and the preparation decoder returns positive blocks
of total trace one. Formula~\eqref{eq:coherentfamily70} therefore gives
exact complete positivity and trace preservation, including on singular
inputs. The chart has at most $4(2B_U+1)^3$ words, which gives the
stated length after adding the inherited preparation payload.
\end{proof}
The reference encoder replays both deterministic encoders when verifying
a supplied target. Bare decoding checks only syntax, rational legality and
the declared construction budget. It cannot certify proximity to an
unspecified target. In particular, reference stability concerns the quantum
instrument defined by the decoded Choi matrices, not access to a classical
processor with equivalent action on an unknown reference-entangled state.
